% GENERATED FILE. Edit canonical lessons, then run npm run generate:books.
% canonical-source-hash: fnv1a64:26223371e83c1978
% canonical-lessons: PA-C202-identity-return, PA-C202-yatri, PA-C202-yatri-review, PA-C202-traveller-label, PA-C202-payoff

\chapter{This Person Is a Traveller}
\label{ch:pa-a2-yatri}
\hlchaptermodality{202}

\begin{chapteropening}
\textbf{By the end of this chapter:} \emph{I can recognise and use the person noun for a traveller in one bounded identity frame.}

Recognise and produce the traveller noun in one bounded identity frame through listening, speaking, reading, and writing.
\end{chapteropening}

\section[\texorpdfstring{\pa{ਇਹ} · \pa{ਹੈ} (eh; hai)}{ਇਹ · ਹੈ (eh; hai)}]{Return to this person and is}

\label{lesson:PA-C202-identity-return}



\begin{quote}
\emph{Hear:} \emph{eh; hai}; YOU SAY: \textbf{\pa{ਇਹ}; \pa{ਹੈ}}
\end{quote}

\begin{grammarlens}[title={Grammar Lens}]
Keep the familiar pointer for a person first and the familiar verb last. The middle slot stays empty until the next lesson owns its new person word.
\end{grammarlens}

\subsection*{Your turn}
\begin{itemize}
\raggedright
  \item \emph{Hear:} \emph{eh}; YOU SAY: \textbf{\pa{ਇਹ}}
  \item \emph{Hear:} \emph{hai}; YOU SAY: \textbf{\pa{ਹੈ}}
  \item Say the first word, pause for the empty middle, then say the last.
\end{itemize}

\begin{tcolorbox}[breakable,colback=teal!4,colframe=teal!35!black,title={Before you move on}]
Produce the two familiar words around the empty middle slot. This opener owns no new word.
\end{tcolorbox}
\section[\texorpdfstring{\pa{ਯਾਤਰੀ} (yaatrii)}{ਯਾਤਰੀ (yaatrii)}]{One person noun for a traveller}

\label{lesson:PA-C202-yatri}



\begin{quote}
Retrieve the familiar first and last words of the identity frame.
\end{quote}

\begin{grammarlens}[title={Grammar Lens: direct owner}]
\begin{quote}
\textbf{\pa{ਯਾਤਰੀ}} — \emph{yātrī} — traveller; person who travels
\end{quote}

Hear and say the whole word before reading it. Keep it as one whole person noun rather than inventing smaller Punjabi pieces. This lesson owns the word alone; the complete identity frame comes only after a separate retrieval.
\end{grammarlens}

\begin{cousinweb}
Current Punjab School Education Board Class 11 material gives this word as the meaning of another traveller term. PSEB's Class 5 book independently uses it for a person who travels beyond Earth. Together they support the bounded person-who-travels meaning in current educational Punjabi. Sources: \href{https://static.pseb.ac.in/media/1749801072\_11\%20\%20Adhunik\%20Punjabi\%202025.pdf}{PSEB Class 11} and \href{https://static.pseb.ac.in/media/1687845902\_Punjabi\%20Pustak-5\%28Pehli\%20Bhasha\%29.pdf}{PSEB Class 5}.

The exact origin of modern written \textbf{\pa{ਯਾਤਰੀ}} remains unknown in the cited sources. Its origin remains unknown here. This lesson therefore makes no etymological claim. Keep the present-day noun whole rather than turning a spelling resemblance into a derivation.
\end{cousinweb}

\subsection*{Your turn}
\begin{itemize}
\raggedright
  \item \emph{Hear:} \emph{yātrī}; YOU SAY: \textbf{\pa{ਯਾਤਰੀ}}
  \item \emph{Read it:} \textbf{\pa{ਯਾਤਰੀ}}
  \item \emph{Write it:} copy \textbf{\pa{ਯਾਤਰੀ}} once, checking the whole word
\end{itemize}

\begin{tcolorbox}[breakable,colback=teal!4,colframe=teal!35!black,title={Before you move on}]
Produce the person noun for “traveller” without looking.
\end{tcolorbox}
\section[\texorpdfstring{\pa{ਯਾਤਰੀ} (yaatrii)}{ਯਾਤਰੀ (yaatrii)}]{Retrieve the new person word on its own}

\label{lesson:PA-C202-yatri-review}



\begin{quote}
\emph{Hear:} \emph{yātrī}; YOU SAY: \textbf{\pa{ਯਾਤਰੀ}}
\end{quote}

\subsection*{Your turn}
\begin{itemize}
\raggedright
  \item \emph{Hear:} \emph{traveller or person who travels}; YOU SAY: \textbf{\pa{ਯਾਤਰੀ}}
  \item \emph{Read it:} \textbf{\pa{ਯਾਤਰੀ}}
  \item \emph{Write it:} cover and write \textbf{\pa{ਯਾਤਰੀ}}
\end{itemize}

\begin{grammarlens}[title={Grammar Lens}]
Keep one person noun and the bounded traveller meaning here. Wider forms, combinations, and nearby transport meanings stay outside this lesson.
\end{grammarlens}

\begin{tcolorbox}[breakable,colback=teal!4,colframe=teal!35!black,title={Before you move on}]
Say and then write the word once from memory.
\end{tcolorbox}
\section[\texorpdfstring{\pa{ਇਹ} \pa{ਯਾਤਰੀ} \pa{ਹੈ} (eh yaatrii hai)}{ਇਹ ਯਾਤਰੀ ਹੈ (eh yaatrii hai)}]{Put the new person noun into one identity frame}

\label{lesson:PA-C202-traveller-label}



\begin{quote}
Retrieve \textbf{\pa{ਇਹ}}, then the new person noun, then \textbf{\pa{ਹੈ}}.
\end{quote}

\begin{grammarlens}[title={Grammar Lens}]
\begin{quote}
\textbf{\pa{ਇਹ} \pa{ਯਾਤਰੀ} \pa{ਹੈ}} — \emph{eh yātrī hai} — this person is a traveller
\end{quote}

Keep the new whole noun between familiar \textbf{\pa{ਇਹ}} and \textbf{\pa{ਹੈ}}. Familiar \textbf{\pa{ਹੈ}} stays last in the already-owned identity frame. This chapter teaches only this one person label; it does not add a form rule or a wider construction.
\end{grammarlens}

\subsection*{Your turn}
\begin{itemize}
\raggedright
  \item \emph{Hear:} \emph{eh yātrī hai}; YOU SAY: \textbf{\pa{ਇਹ} \pa{ਯਾਤਰੀ} \pa{ਹੈ}}
  \item \emph{Read it:} \textbf{\pa{ਇਹ} \pa{ਯਾਤਰੀ} \pa{ਹੈ}}
  \item \emph{Write it:} cover and write \textbf{\pa{ਇਹ} \pa{ਯਾਤਰੀ} \pa{ਹੈ}}
\end{itemize}

\begin{tcolorbox}[breakable,colback=teal!4,colframe=teal!35!black,title={Before you move on}]
Produce the complete person identity frame.
\end{tcolorbox}
\section[\texorpdfstring{\pa{ਇਹ} \pa{ਯਾਤਰੀ} \pa{ਹੈ} (eh yaatrii hai)}{ਇਹ ਯਾਤਰੀ ਹੈ (eh yaatrii hai)}]{A bounded traveller-noun payoff}

\label{lesson:PA-C202-payoff}



\begin{quote}
Retrieve the new person noun alone before rebuilding its identity frame.
\end{quote}

\subsection*{Your turn: four skills}
\begin{itemize}
\raggedright
  \item \emph{Hear:} \emph{eh yātrī hai}; YOU SAY: \textbf{\pa{ਇਹ} \pa{ਯਾਤਰੀ} \pa{ਹੈ}}
  \item \emph{Read it:} \textbf{\pa{ਇਹ} \pa{ਯਾਤਰੀ} \pa{ਹੈ}}
  \item \emph{Write it:} cover the line and write \textbf{\pa{ਇਹ} \pa{ਯਾਤਰੀ} \pa{ਹੈ}}
\end{itemize}

\begin{tcolorbox}[breakable,colback=teal!4,colframe=teal!35!black,title={Before you move on}]
Say the new noun, then produce the complete identity frame.

This is the twenty-fifth prerequisite tranche exposed by Punjabi A2 mock drafting. It does not unblock Punjabi A2 mock 1 by itself. It is not an A2 mock, not A2 readiness or attainment, and not external or human validation.
\end{tcolorbox}
